\documentclass[11pt]{article}
\usepackage[margin=1in]{geometry}
\usepackage{amsmath,amssymb,amsthm,mathtools,booktabs,longtable,hyperref}
\newtheorem{theorem}{Theorem}[section]
\newtheorem{lemma}[theorem]{Lemma}
\newtheorem{proposition}[theorem]{Proposition}
\newtheorem{corollary}[theorem]{Corollary}
\theoremstyle{remark}\newtheorem{remark}[theorem]{Remark}
\newcommand{\dd}{\,\mathrm d}
\newcommand{\E}{\mathbb E}
\newcommand{\one}{\mathbf 1}
\newcommand{\eps}{\varepsilon}
\newcommand{\Rhalf}{\mathcal R}
\title{The late $3\pi/2$ strip: delay, chart contraction and an arrival-weighted budget}
\author{Analytical increment AN03\_late\_strip\_budget\_v1}
\date{October 26, 2026}
\begin{document}\maketitle

\section{Scope, sources and conclusion}
This increment proves PISB's remaining arrival-weighted strip consumer on
the stopped interval of S2CP/CFA, retaining their non-strip inputs only.
In particular the fixed-chart field allowance, core/reference geometry,
duration, residual guards and the improved secondary bound are retained
exactly as in S2CP. Neither CFA's strip input nor its passage conclusion is
used. No partition is enlarged and no weak-passage bootstrap is executed.
Source abbreviations and exact repository provenance are recorded in the
companion review. An external reference such as MOMC \texttt{jacobian}
means the statement with that suffix in the cited source, not a new
assumption. All new labels have prefix \texttt{an03lsbv1:}.

Explicitly, after $t_0$ the retained guards include
$0\le c\le c_b<1$, $t_*-t_0\le L_T\ell$, $\int r\le B$,
$J_G<q^2K_Y^2$, $E_+<E_g$ and
$\|e_2\|_{P_1}+\|e_1\|_{P_2}\le M_{\rm res}q$.
The instantaneous core chart and mass bound, reference offset $e<e_*$,
and S2CP's instantaneous outside guard
$\sup\mu_{\mathcal B}/q=b_B(q)=o(1)$ are retained as stopped guards.
The last guard is not an assumed time-integrated strip budget.
The inherited c-form has positive donor weights, donor mass at most
one, and $\int\epsilon_F\le\overline F_F$ in value and induced
$\ell^1$ receiver derivative on the fixed chart
$[-L_c,L_c]^2$, where $L_c=R_c+1$ and $R_c=79/51$.
S2CP improves its temporary secondary guard under these inputs.
The current result supplies a quantitative strip budget under these
same guards; it does not assert strict improvement of an arbitrarily
prescribed function $b_B(q)$ for the full outside population.

Put $\ell=\log(1/q)$ and use the following notation, keeping the noise
correction in $\beta$:
\begin{align}
 &\lambda=\rho^2\in[169/400,49/100],\quad Q=\lambda+q^2,
 \quad\beta=\lambda/Q,\quad \omega=(1-\lambda)/2,\nonumber\\
 &a=\rho K(\rho a),\quad P_\rho=\Phi(\rho a),\quad
 \alpha=\lambda P_\rho/2,\quad d=1/2-\alpha,
 \quad s_{\rm pr}=d/\omega,\quad p=3/s_{\rm pr},\nonumber\\
 &c_m=4/(1-\lambda)+1/2,\quad u=10-c_m,\quad v=u-1,
 \quad k=10(1-\lambda),\quad t_m=c_m\ell,\nonumber\\
 &M_0=A_q q^u,\quad c_A\le A_q\le C_A,\qquad
 t_0=10\ell-\log A_q+O(q^2\ell^2+q^u).
 \label{an03lsbv1:parameters}
\end{align}
Here $K(z)=\phi(z)+z\Phi(z)$ is the Gaussian positive-part kernel;
$M$ is MOMC's original right-half radial mass, and $M(t_0)=1/2$.
Fix $0<\nu\le1/100$ and take $q$ small enough that
$1/\beta=1+q^2/\lambda\le1+\nu$. The symbol $r$ below, when used in
a stopped guard, is $\sup_\theta|R_{11}(\theta)|$, not a receiver angle.

Write $b_0=-\pi/2$ using the continuous lift at $3\pi/2$, and set
\begin{equation}
 J_3=\{b=b_0-q s_0:0\le s_0<s_q\},\qquad
 s_q^2=\frac{1+\eps}{k_G}\log\ell,\quad
 k_G=169/3200,\quad 0<\eps<1/2.
 \label{an03lsbv1:strip}
\end{equation}
Endpoint inclusion is immaterial for mass but is retained in every
labelwise assertion. Initial mass is $m(0,b)=q^{10}$ and
$\dd\lambda_0=\dd b/(2\pi)$ in absolute orientation.

The stopped clock and cap-one conclusions of S2CP and PISB give
\begin{align}
 I(t)&=Q\int_{t_0}^t(1-c),\qquad
 e^{I(t)}\le F_q q^{-k},\qquad
 F_q=C_H E_g A_q^\lambda e^{D_{\rm clk}},
 \label{an03lsbv1:clock}\\
 D_{\rm clk}&=O(q^{1/8}+q^2\ell),\qquad
 m(t,b)\le C_{\rm abs}m(t_0,b)e^{I(t)},\nonumber\\
 C_{\rm abs}&=e^{2B+\eta_q},\quad
 \eta_q=C(1+M_{\rm res})qL_T\ell=o(1).
 \label{an03lsbv1:cap}
\end{align}
The constants $C_H,c_A,C_A$ are the source constants. The duration is
$t_*-t_0\le L_T\ell$. Fix S2CP's $C_D,A_0$ and $C_g=e^{6B+3}$;
$A_0=C_D^2\max(1,\delta_0^{-2})$. Increasing $C_g$ by a fixed
model factor, if required by equivalent norm conventions, is harmless.

Let $\tau_b$ be the first joint arrival after $t_0$ in the fixed S2 chart,
with $\tau_b=t_0$ for receivers already there and $\tau_b=\infty$ if
there is no arrival. Let $\xi_b$ be the joint scaled displacement from
the actual same-field reference at arrival. Use the same norm as S2CP;
its equivalence to the maximum norm costs a fixed factor. Define
\begin{align}
 W_3(t)={}&\int_{J_3\cap\{\tau_b>t\}}m(t,b)\dd\lambda_0(b)\nonumber\\
 &+C_g\int_{J_3\cap\{\tau_b\le t\}}m(\tau_b,b)
 \max\{1,A_0q^2\xi_b^2e^{\beta[I(t)-I(\tau_b)]}\}^{1+\nu}
 \dd\lambda_0(b).
 \label{an03lsbv1:consumer}
\end{align}

\begin{theorem}[The remaining strip budget]
\label{an03lsbv1:main}
On the stopped interval just specified, the gain in
\eqref{an03lsbv1:consumer} is arrival-valid and $\mu_{J_3}(t)\le W_3(t)$.
Define
\begin{align}
 \chi_{\rm lag}&=1/4,\qquad a_e=v+du,\qquad
 H_q=C A_0 A_q^{2d}F_q^\beta,\nonumber\\
 \delta_B&=c_m-(1+\nu)k,\qquad
 \delta_L=\chi_{\rm lag}-\nu k,\nonumber\\
 \delta_E&=2+(1+\nu)(2-\beta k+12d),\qquad
 \delta_F=2+(1+\nu)(2-\beta k+2a_e).
 \label{an03lsbv1:exponents}
\end{align}
There is a fixed model/S1 constant $C$ such that, uniformly in $t$,
\begin{align}
 \frac{W_3(t)}q\le C A_q^{-1} C_g C_{\rm abs}
 \big\{&F_q^{1+\nu}(s_q q^{\delta_B}+q^{\delta_L})+q^2\nonumber\\
       &+H_q^{1+\nu}(q^{\delta_E}+q^{\delta_F})\big\}.
 \label{an03lsbv1:budget}
\end{align}
The constants in $H_q$ and the displayed $C$ do not depend on LCRC's
guards. All guard dependence is displayed in the fixed source gain and
clock factors. In particular,
\begin{equation}
 \sup_{[t_0,t_*]}\frac{W_3}q
 \le C_{\rm fixed}(1+s_q)q^{769/4000}=o(1),\qquad
 \int_{t_0}^{t_*}\frac{W_3}q
 \le L_T\ell C_{\rm fixed}(1+s_q)q^{769/4000}=o(1).
 \label{an03lsbv1:finalbudget}
\end{equation}
Together with PISB's $\pi$ strip, middle and outer zones this supplies
CFA's whole-strip input with both budgets vanishing. This is an
improvement of that input on the retained stopped trajectory; it is
not an application of CFA's passage theorem.
\end{theorem}

\section{Target-only delay, with the lower endpoint regularized}
\begin{lemma}[Sign, exact time shift and delay density]
\label{an03lsbv1:delay}
For small $q$, the aligned target velocity is positive on
$[b_0-qs_q,b_0]\cup(b_0,qa_q)$, where $a_q$ is PROFILE's exact root.
For $b\in J_3$ put
\begin{equation}
 T(b)=\int_b^{b_0}\frac{\dd\vartheta}{v_q(\vartheta)}.
 \label{an03lsbv1:delaydef}
\end{equation}
Then, writing a star for the endpoint $b_0$ trajectory,
\begin{align}
 \widehat\theta(t,b)&=\widehat\theta_*(t-T),\qquad t\ge T,\nonumber\\
 \widehat m(t,b)&=L(b)\widehat m_*(t-T),\qquad
 L(b)=\exp\left(2\int_0^T g_q(\widehat\theta(w,b))\dd w\right)
       =e^{o(1)},\quad L(b)\ge1,\nonumber\\
 \widehat\theta(t,b)&\le\widehat\theta_*(t),\qquad
 \widehat m(t,b)\le\widehat m_*(t)\quad(t\ge0),
 \label{an03lsbv1:shift}\\
 \dd\lambda_0&=\frac{v_q(b(T))}{2\pi}\dd T
 \le CqK(-s_0(T))\dd T\le Cq\dd T.
 \label{an03lsbv1:density}
\end{align}
The first identity also holds for $0\le t<T$ if the endpoint orbit
is continued backwards to time $-T$. At $b=b_0$, $T=0$ and $L=1$.
\end{lemma}
\begin{proof}
PROFILE's exact equation and the target radial equation are
\[
 v_q(\theta)=\frac q2\{-\sin\theta K(\cos\theta/q)
             +\rho\cos\theta K(\rho\sin\theta/q)\},\quad
 (\log\widehat m)'=2g_q(\theta),\quad g_q'=2v_q.
\]
For $\theta=-r\in(-\pi/2,0)$,
\[
 v_q(-r)=\tfrac12\sin r\cos r+	frac q2\sin r K(-\cos r/q)
              +\tfrac{\rho q}2\cos r K(-\rho\sin r/q)>0.
\]
The positive sign up to $qa_q$ follows from PROFILE's root and sign
lemma (on the remaining fixed scaled interval). In the initial layer,
\[
 v_q(b_0-qs)=\frac q2\left\{\cos(qs)K(-\sin(qs)/q)
                 -\rho\sin(qs)K(-\rho\cos(qs)/q)\right\}.
\]
Uniformly for $0\le s\le s_q$, this is comparable to $qK(-s)>0$.
Indeed $K(-s)\ge c(1+s)^{-2}e^{-s^2/2}$, by integrating its
Gaussian integral over an interval of length $(1+s)^{-1}$; the second
term is exponentially small in $q^{-2}$. The first term changes by a
relative $O(q^2(1+s)^4)=o(1)$ under its trigonometric replacements.
Consequently
\[
 T_{\max}\le C(1+s_q)^3e^{s_q^2/2},\qquad q^2T_{\max}=o(1).
\]
Both projected cluster means are nonpositive in this initial layer,
so $0\le g_q\le Cq^2$ there. Autonomy gives the time shift and exact
radial prefactor, whose logarithm is at most $Cq^2T_{\max}=o(1)$.
Scalar order and $g_q'=2v_q>0$ along the entire relevant interval
give equal-time mass domination, including times before emergence.
Finally differentiate \eqref{an03lsbv1:delaydef} in absolute orientation
to obtain \eqref{an03lsbv1:density}. No zero-seed convention is used.
\end{proof}

\begin{lemma}[The endpoint's seed, arrival and uniform local decay]
\label{an03lsbv1:endpoint}
On logarithmic time intervals the endpoint obeys
\begin{align}
 \widehat m_*(t)&\le C(q^{10}+q^{12}e^t),
 \label{an03lsbv1:massendpoint}\\
 |\widehat\theta_*(t)|&\le Cq^{-1}e^{-t/2}quad
       \text{if }t\ge2\ell+C\text{ and }\widehat\theta_*(t)<0.
 \label{an03lsbv1:angleendpoint}
\end{align}
The times to $-\pi/4$ and to $-qL$, for fixed sufficiently large $L$,
are $2\ell+O(1)$ and $4\ell+O(1)$, respectively. After a further fixed
time the trajectory is in the interior of $[-R,R]$, $R=79/51$.
Uniformly after that arrival and for $q^2t=o(1)$,
\begin{equation}
 \widehat m_*(t)\asymp q^{12}e^t,\qquad
 a_q-\widehat\theta_*(t)/q\asymp q^{-4d}e^{-dt}.
 \label{an03lsbv1:endpointprofile}
\end{equation}
All constants are model/S1 constants, uniform on the ratio box.
\end{lemma}
\begin{proof}
The target coordinate equation at $(\widehat U_1,\widehat U_2)
=(0,-q^5)$ gives $\widehat U_1(1)\asymp q^6$, a strictly positive
Gaussian seed. As long as $\widehat U_2<0$ its negative part is at
most $q^5e^{q^2t}$ (MOMC's coordinate estimate). The $P_1$ Stein
identity gives $\widehat U_1'\ge\widehat U_1/2$ once it is positive.
For an upper bound the same coordinate equation gives
\[
 \widehat U_1'\le(1/2+Cq+q^2)\widehat U_1
                       +Cq^6e^{q^2t}.
\]
Thus $\widehat U_1\le Cq^6e^{t/2+Cqt}$ before the zero crossing of
$\widehat U_2$, proving the two upper bounds there. The zero crossing
occurs at $4\ell+O(1)$ with mass $\asymp q^8$, by the fixed-chart
arrival calculation below. Thereafter the target spectral bound
$(\log\widehat m)'\le1+2q^2$ proves
\eqref{an03lsbv1:massendpoint} for the rest of the interval.

Use PROFILE's Q4 reciprocal-speed calculation, retaining its
$\cos b+q$ regularization. Its arrival time is
$2\log(1/[q(\cos b+q)])+O(1)$; at $b=b_0$ this is $4\ell+O(1)$.
Stopping that same calculation at $-\pi/4$ gives $2\ell+O(1)$.
On any subsequent fixed scaled interval the target scalar field has
its exact root $a_q$, derivative $-d+O(q^2)$ there, and a quadratic
Taylor remainder bounded by $C|x-a_q|^2$. After a fixed further time,
attraction gives $\int|x-a_q|\le C$, independently of the terminal
time. Integration of the logarithmic displacement equation gives
bounded distortion $e^{O(1)+O(q^2t)}$. The local radial error has the
same integrable nonlinear contribution and $O(q^2t)$ correction.
Matching at $4\ell+O(1)$ proves \eqref{an03lsbv1:endpointprofile}.
This proves the needed uniform-in-time Q4 specialization; it does
not apply PROFILE's common, later Q1 time guard at an earlier time.
\end{proof}

\section{Transfer and a chart lemma for an arbitrary receiver}
Q2EC's all-label comparison at $t_m$ is applicable to both $b$ and
$b_0$ because the original total mass is still $O(q^u)$. It gives
\begin{equation}
 |\theta-\widehat\theta|+|\psi-\widehat\theta|\le Cq^u,
 \qquad m(t_m,b)=\widehat m(t_m,b)e^{O(q^u)}.
 \label{an03lsbv1:early}
\end{equation}
Fix the following label sets by their states at this single time:
the charted set has both scaled angles in $[-R,R]$; of its complement,
the boundary set has $\theta(t_m)\le b_0+2q$, and the transit set is
the rest. Every non-charted member has $\sin\theta(t_m)<0$ for small
$q$: if $\theta\ge0$, target order and \eqref{an03lsbv1:early} put both
angles between $-o(q)$ and $q(a_q+o(1))$, strictly in the chart.
These are fixed ancestry subsets, not time-dependent masks.

The target delay also defines fixed label subsets and may be used
simultaneously with this state classification. Equations
\eqref{an03lsbv1:shift}--\eqref{an03lsbv1:massendpoint}, including the
case $T>t_m$, imply
\[
 m(t_m,b)\le C\{q^{10}+q^{12}e^{t_m-T(b)}\}.
\]
PISB's all-receiver bridge uses $\int_{t_m}^{t_0}S\le3/2+o(1)$, and
therefore yields for every $b\in J_3$
\begin{equation}
 m(t_0,b)\le C A_q^{-1}\{q^{c_m}+q^2e^{-T(b)}\}.
 \label{an03lsbv1:entrymass}
\end{equation}
The fixed bridge cost is $e^{6+o(1)}$. This bound does not presume
that the receiver has arrived by $t_0$.

\begin{lemma}[Single-label contraction in the strong chart]
\label{an03lsbv1:chart}
An arbitrary original receiver that is in $[-R,R]^2$ at time
$s\in[t_m,t_0]$ remains in that chart and satisfies, relative to
MOMC's actual same-field reference,
\begin{equation}
 \xi(t_0)\le C_{\rm sh}\xi(s)G(1/2,M(s))
       \le C\xi(s)(2M(s))^d,
 \quad
 G(M_1,M_0)=(M_1/M_0)^{-d}
       \left(\frac{1-M_0}{1-M_1}\right)^\alpha.
 \label{an03lsbv1:single}
\end{equation}
The constants depend only on the model and the fixed chart, not on
LCRC guards or the receiver's ancestry. The exact finite-noise field
need not be cooperative.
\end{lemma}
\begin{proof}
MOMC's reduction holds for any receiver in a fixed chart. Its leading
Jacobian, with the actual right-half donors held fixed, is
\[
 Df(x,y)=\begin{pmatrix}
 -(1-M)/2+\rho^3(y-x-\Pi(x))\phi(\rho x)/2&\lambda\Phi(\rho x)/2\\
 \lambda\Phi(\rho x)/2&-(1-M)/2
 \end{pmatrix}.
\]
Here $\Pi(x)=\int(\xi-x)_+\dd\mu_s\ge0$.
The second row sum is at most $-1/4+\lambda/2\le-1/200$.
For the first row drop $\Pi$ and replace $y$ by $R$. The remaining
positive expression is
\[
 H(x)=\lambda\Phi(\rho x)/2+\rho^3(R-x)\phi(\rho x)/2.
\]
Since $H'(x)=-\rho^5x(R-x)\phi(\rho x)/2$, its maximum on
$[-R,R]$ is at zero. Using $\phi(0)<2/5$,
\[
 H(0)\le\frac{49}{400}+\frac{27097}{255000}
        <\frac{49}{200}.
\]
Thus the induced maximum-norm logarithmic norm of $Df$ is at most
$-1/200$ throughout the square, not only near the atomic center.
MOMC's error is bounded by $C_Rq^2+160\mu/q=o(1)$ there,
since $\mu\le Cq^2\ell^2M$. The exact logarithmic norm is consequently
at most $-1/400$. Absolute off-diagonal values change by at most the
$C^1$ error; finite-noise cooperativity is unnecessary.

For completeness the square is invariant. Replace the coherent donor
state by $(a,a)$, with error $o(1)$ uniformly, using MOMC's mean and
spread conclusions. On the upper corner the two leading components
equal $f_0(R)+M(R-a)/2$, where
$f_0(x)=(-x+\rho K(\rho x))/2$ and $f_0(a)=0$.
Its derivative is at most $-1/200$ for $M\le1/2$.
At the lower corner the second component is at least
$(R+a)/200$, and the first is no smaller: for $x<a$,
$\rho F(\rho x,\rho a)\le a$, where $F$ is MOMC's donor kernel.
Positive leading off-diagonal entries reduce the face checks to
these corners. The resulting fixed strict inward margins survive
the $o(1)$ perturbation. This proves invariance and exponential
contraction at rate $1/400$ to the actual reference, hence
$\int_s^{t_0}\xi\le C_R$.

Now retain the sharper central Jacobian
$J_{\rm sh}(M)$, whose row sum is $-d+M/2$. On the segment from
the actual reference to the receiver its difference from the exact
Jacobian has norm at most
\[
 C_R\{\xi+e+M r_{\Rhalf}+q^2+\mu/q\}.
\]
Here $e$ is the reference offset from $(a,a)$ and $r_{\Rhalf}$ the
right-half donor spread. MOMC gives vanishing integrals for all
terms except $\xi$, whose integral was just bounded independently
of the interval length. Logarithmic-norm integration therefore gives
$\xi(t_0)\le C\xi(s)\exp\int_s^{t_0}(-d+M/2)$.
MOMC's mass equation, with integrated defect $O(q^2\ell^2)$, converts
this exponential to $G$ up to a fixed factor. Since $M(t_0)=1/2$,
the last factor $(2(1-M(s)))^\alpha$ is bounded by $2^\alpha$.
\end{proof}

\begin{remark}
MOMC's small-displacement pair proof does not apply verbatim to a
new receiver at the edge of the chart. The uniform logarithmic-norm
calculation above supplies the missing bounded-distortion argument.
It does not assert a multiplicative $e^{o(1)}$ shape comparison for
this new ancestry.
\end{remark}

\section{A delay class that arrives before half mass}
Choose the fixed delay cutoff
\begin{equation}
 T_c=(k-2+\chi_{\rm lag})\ell=(k-7/4)\ell,\qquad \chi_{\rm lag}=1/4,
 \quad J_E=\{T\le T_c\},\quad J_L=\{T>T_c\}.
 \label{an03lsbv1:split}
\end{equation}
The delay parameter $\chi_{\rm lag}=1/4$ is unrelated to TAILREP's
core-cutoff exponent $\chi=1/20$, which is unchanged.
These are fixed label sets. No arrival assertion for $J_L$ will be
needed to pay its budget.

\begin{lemma}[Transit and entry for the early delay class]
\label{an03lsbv1:transit}
Every $b\in J_E$ enters $[-R,R]^2$ by
\begin{equation}
 s_b\le\max\{t_m,4\ell+T(b)+C\}<t_0.
 \label{an03lsbv1:arrivaltime}
\end{equation}
At $t_0$ it obeys the mass bound \eqref{an03lsbv1:entrymass} and
\begin{equation}
 \xi(t_0,b)\le C A_q^d
       \{q^{6d}e^{dT(b)}+q^{a_e}\},\qquad a_e=v+du.
 \label{an03lsbv1:massshape}
\end{equation}
\end{lemma}
\begin{proof}
PISB's rational bound gives
$c_m-k\ge15259/9240>33/20$. Thus at $t_m$ every member of $J_E$
has target age at least $(c_m-k+7/4)\ell$, greater than
$(17/5)\ell$. By \eqref{an03lsbv1:angleendpoint} and
\eqref{an03lsbv1:early}, its negative angles have magnitude
$O(q^{7/10})$; its positive angles are at most $q(a_q+o(1))$.
Also $\Delta(t_m):=\theta(t_m)-\psi(t_m)=O(q^u)$.

Here are the receiver estimates needed between this state and a
fixed chart. Stop temporarily in a small fixed cone around zero.
The $P_1$ receiver gate is exponentially full. With MOMC's strong
donors and arbitrary complement of mass $\mu$,
\begin{align}
 R_{11}&=(1-M)/2+O(q^2+\mu),\nonumber\\
 R_{22}&\ge-C(q^2+\mu),\qquad
 |R_{12}-R_{21}|\le C(qM+\mu),\nonumber\\
 |R_{12}|+|R_{21}|&\le C(q+\mu),\qquad
 |R_{22}\sin\theta|\le C(q+\mu)\quad(\theta\le0).
 \label{an03lsbv1:transitmatrix}
\end{align}
To verify every estimate, remove the $P_1$ gates and use
$\E_{P_1}XX^T=\operatorname{diag}(1+q^2,q^2)$;
the donors have $W_1U_1=m(1+O(q^2))$ and $W_2=O(q\sqrt m)$.
On $P_2$, $f_{\Rhalf,1}$ has norm $O(qM)$ and
$f_{\Rhalf,2}$ norm $O(q^2M)$, since all donor inputs and outputs
are in the strong chart. These give respectively $O(q^2M)$ in
$R_{11}$, $O(q^2M)$ in $R_{22}$ and $O(qM)$ in the cross entries.
The target selected covariance is symmetric and positive
semidefinite, so its contribution cancels from $R_{12}-R_{21}$
and is nonnegative in $R_{22}$.
For $\theta\le0$ the projected $P_2$ margin gives
$|\sin\theta|\E_{P_2}[X_2^2\one_{a_\theta\cdot X>0}]
\le Cq$; Gaussian conditioning proves this also at $\theta=0$.
PROFILE's unscaled outside value bound is $C\mu$ in each entry.
These observations prove \eqref{an03lsbv1:transitmatrix} without
extending the fixed-chart reduction to an unbounded chart.

The exact difference identity is
\[
 \Delta'=-(\operatorname{tr}R_\theta)\sin\Delta
                 +(R_{12}-R_{21})\cos\Delta.
\]
The trace is at least $1/5$ for small $q$ before half mass.
MOMC's increasing mass and $\mu\le Cq^2\ell^2M$ imply
\begin{equation}
 |\Delta(t)|\le Cq^u e^{-(t-t_m)/6}+CqM(t).
 \label{an03lsbv1:delta}
\end{equation}
The cone and the small-$\Delta$ restrictions are preserved by the
following inequalities, so they are legitimate first-exit devices.
For $z=-\theta\ge0$, the exact angular equation and
\eqref{an03lsbv1:transitmatrix} give
\[
 z'\le-\tfrac12(1-M)z+Cz^3+Cq.
\]
Here replacement of $\psi$ by $\theta$ costs $C|\Delta|=O(q)$;
the initial $q^u$ term is smaller than $q$.
Fix a sufficiently large model constant $L$. Whenever $z\ge qL$,
smallness of the cone and largeness of $L$ give $z'\le-z/8$.
Consequently $\int z^2\le C z(t_m)^2$ and
$\int q/z\le8/L$ up to that hitting time (change variables using
$-z'\ge z/8$). Integrating the sharper logarithmic inequality gives
\[
 z(t)\le C z(t_m)\exp\{-\tfrac12(t-t_m)+\tfrac12\int_{t_m}^tM\}.
\]
Since $\int M\le3/2$, \eqref{an03lsbv1:angleendpoint} puts the hit
of $-qL$ no later than $4\ell+T+C$. Receivers already above $-qL$
need only the following bounded-time argument.

For a non-charted label the target attraction and
\eqref{an03lsbv1:early} imply $t_m\le4\ell+T+C$; otherwise it
would already lie strictly inside $[-R,R]^2$. Thus all nontrivial
arrivals considered here occur by
\[
 (k+2+\chi_{\rm lag})\ell+C\le(321/40)\ell+C,
 \qquad M\le C A_q q^{79/40}=o(1),
\]
using $M(t)\le C A_q q^{10}e^t$, a direct consequence of MOMC's
mass law. Equation \eqref{an03lsbv1:delta} then gives
$\Delta/q=o(1)$. On the fixed enlarged chart containing
$[-L,L]^2$, MOMC's reduction is uniformly the target field plus
$o(1)$ on this bounded additional time interval. The aligned target
orbit from $[-L,0]$ enters a fixed interior subinterval of $[-R,R]$
in bounded time. Finite-time continuous dependence, with the
displayed $o(1)$ misalignment, proves joint entry in $[-R,R]^2$.
The same argument covers the bounded transition band at $t_m$.
This closes the cone/time first exit and proves
\eqref{an03lsbv1:arrivaltime}; the gap to $t_0$ tends to infinity.

If the target label is already a fixed distance inside the chart at
$t_m$, \eqref{an03lsbv1:endpointprofile}, the absolute comparison
and the actual reference offset $O(q^2+q^v)$ give
\[
 \xi(t_m,b)\le C\{q^{d(c_m-4)}e^{dT}+q^v\}.
\]
Here $q^2=O(q^v)$ uniformly on the box. Applying
Lemma~\ref{an03lsbv1:chart} multiplies this by at most
$C A_q^d q^{du}$ and proves \eqref{an03lsbv1:massshape}.
For the remaining arrivals, $\xi(s_b)\le C$ and
$M(s_b)\le C A_q q^6e^T$, so the same lemma gives the first
term of \eqref{an03lsbv1:massshape}. This includes the bounded
transition band and proves the result without resetting the reference.
\end{proof}

\section{Restarting the gain and paying the weighted consumer}
\begin{lemma}[Gain from the actual arrival]
\label{an03lsbv1:restart}
For every finite $\tau_b\le t_*$ in \eqref{an03lsbv1:consumer},
\begin{equation}
 m(t,b)\le C_gm(\tau_b,b)
 \max\{1,A_0q^2\xi_b^2e^{\beta[I(t)-I(\tau_b)]}\}^{1+\nu}
 \quad(t\ge\tau_b).
 \label{an03lsbv1:arrivalgain}
\end{equation}
\end{lemma}
\begin{proof}
This is a restart of S2CP's distance and DECG's gain argument,
not a new entry-profile assertion. At arrival the receiver is on
the strong branch in a fixed chart. The c-form local distance
estimate and GESG's matrix-value splice are estimates from an
arbitrary start time; their constants use only upper bounds on
$\int r$, the nonnegative gate excess, and the field error budget.
Each subinterval inherits the same bounds $B$, $\mathfrak X$ and
$\overline F_F$. The clock difference is $I(t)-I(\tau_b)$, with
the same positive derivative lower bound. Thus the distance
estimate holds with the same $C_D$ until physical exit, including
reentries. The signed radial estimate and the reference radial
budget in DECG then give
\[
 e^{6B+1/\beta+o(1)}
 \max\{1,A_0q^2\xi_b^2e^{\beta[I(t)-I(\tau_b)]}\}^{1/\beta}.
\]
Its small errors are bounded by full-interval small errors, uniformly
over start times and receivers. Since $1/\beta\le1+\nu$, this is
\eqref{an03lsbv1:arrivalgain} with $C_g=e^{6B+3}$ for small $q$.
None of the entry-tail or ancestry assumptions enters this receiver
gain proof. At $\xi_b=0$ uniqueness gives the actual reference;
its radial factor is already included in $C_g$.
\end{proof}

\begin{proof}[Proof of Theorem~\ref{an03lsbv1:main}]
The restart lemma and the unarrived term imply $\mu_{J_3}\le W_3$.
We estimate the consumer itself, rather than only this smaller mass.
All arrival displacements are bounded by a fixed chart constant.
Combining \eqref{an03lsbv1:cap} at arrival with the multiplier in
\eqref{an03lsbv1:consumer} shows, labelwise, that its contribution is
at most
\begin{equation}
 C C_g C_{\rm abs}m(t_0,b)e^{(1+\nu)I(t)}.
 \label{an03lsbv1:coarseconsumer}
\end{equation}
For small $q$, $A_0q^2\xi_b^2\le1$ for all bounded arrival states.
The estimate follows from $I(\tau_b)\le I(t)$ and
$\beta(1+\nu)\le1+\nu$. For an unarrived label it follows directly
from cap one. Hence it does not require that a late label remain
on the lower side, or ever arrive. GESG's lower-side no-gain is
consistent but is not needed for this coarser bound.

The first summand $C A_q^{-1}q^{c_m}$ in
\eqref{an03lsbv1:entrymass}, integrated over the whole strip of
width $qs_q$, costs in $W_3/q$ at most
\[
 C A_q^{-1}C_g C_{\rm abs}F_q^{1+\nu}s_q
                          q^{c_m-(1+\nu)k}.
\]
This is an upper-bound decomposition, not a decomposition of the
actual radial measure into evolving subpopulations.
For the second mass summand restricted to $J_L$, the density
\eqref{an03lsbv1:density} gives
\[
 \int_{J_L}q^2e^{-T}\dd\lambda_0\le Cq^3e^{-T_c}.
\]
Equation \eqref{an03lsbv1:coarseconsumer} therefore pays its $W_3/q$
by $C A_q^{-1}C_g C_{\rm abs}F_q^{1+\nu}q^{\chi_{\rm lag}-\nu k}$.
This explicitly retains both post-$t_0$ arrivals and the never-arrived
population in the estimate.

Every member of $J_E$ is charted at $t_0$. For its correlated mass
summand $C A_q^{-1}q^2e^{-T}$ use
\eqref{an03lsbv1:massshape} and the arrival gain at $t_0$. Up to a
fixed factor the gain integrand is bounded by
\[
 1+\left[H_q q^{2-\beta k}
              \{q^{12d}e^{2dT}+q^{2a_e}\}\right]^{1+\nu}.
\]
After division by $q$, integration with
$\dd\lambda_0\le Cq\dd T$ and $\int_0^\infty e^{-T}\dd T=1$ gives
\[
 C A_q^{-1}C_g q^2\left\{1+
 \frac{H_q^{1+\nu}q^{(1+\nu)(2-\beta k+12d)}}
      {1-2d(1+\nu)}
 +H_q^{1+\nu}q^{(1+\nu)(2-\beta k+2a_e)}\right\}.
\]
The extension of the nonnegative integral to $T=\infty$ is only
an estimate; it does not extend the early-class shape assertion.
Analytically, $d\le631/1600$, since $P_\rho\ge1/2$, so
\begin{equation}
 1-2d(1+\nu)\ge\frac{16269}{80000}>0.
 \label{an03lsbv1:momentmargin}
\end{equation}
This proves \eqref{an03lsbv1:budget}. It is a full weighted moment,
hence also bounds any restriction to TAILREP's cutoff. The correlated
part has delay tail index $1/d$, and
$1/d>3\omega/d=p$ because $3\omega\le693/800<1$.
The extra $q^{a_e}$ floor and the uncorrelated mass summand have
been charged separately; no two-sided raw profile is claimed.

For clarity, the positive powers can all be checked with rational
bounds. From the sources and \eqref{an03lsbv1:parameters},
\[
 d>7/20,\quad u\ge169/102,\quad v\ge67/102,\quad
 a_e\ge841/680,\quad k\le231/40.
\]
Consequently
\begin{align*}
 \delta_B&\ge15259/9240-231/4000>3/2,\\
 \delta_L&\ge1/4-231/4000=769/4000,\\
 \delta_E&\ge2+17/40=97/40,\\
 \delta_F&\ge2-(101/100)(177/136)=9323/13600.
\end{align*}
Here $\beta\le1$ was used only to bound the exponents from below;
it remains exact in \eqref{an03lsbv1:budget}. All other factors are
fixed or tend to fixed limits. Multiplication by $L_T\ell$ proves
\eqref{an03lsbv1:finalbudget}. Adding PISB's already proved zone
budgets proves the final strip assertion. In particular CFA's
small fixed strip allowance is met by choosing $q$ last, without
using its passage conclusion.
\end{proof}

\section{Constants and limits of the conclusion}
The chart radius $R=79/51$, the transit radius $L$, its bounded
additional arrival time, $C_{\rm sh}$, the delay and endpoint
constants, the bridge factor $e^6$, and the bounds $c_A,C_A$ depend
only on the model/S1 sources. The new cutoff $\chi_{\rm lag}=1/4$ is fixed
numerically by an analytic margin, independently of any guard.
The middle-zone parameter $\eps$ is inherited from PISB. None of
these quantities enlarges a LCRC core, entry, or reference constant.

The stopped constants enter through $C_D$, $A_0$, $C_g=e^{6B+3}$,
$C_{\rm abs}$, $L_T$ and $F_q$. Thus the three types of displayed
$A_q$ factors are $A_q^{-1}$, $A_q^{\lambda(1+\nu)-1}$ in the
coarse consumer, and
$A_q^{-1+(2d+\lambda\beta)(1+\nu)}$ in the shape moments.
The additional fixed factors are powers of $E_g$, $C_H$, $A_0$
and $e^{D_{\rm clk}}$ exactly as in \eqref{an03lsbv1:budget}.
There is no hidden choice of a large chart depending on $B$ or
$M_{\rm res}$ in the pre-entry proof. The post-entry chart and gain
constants are the already fixed S2CP ones, and multiply positive
powers of $q$ only.

The source constant order is unchanged: model and S1 geometry;
the retained stopped guards and field allowances with S2CP's acyclic
order (including its later choice of $C_\Sigma$ after $C_0$ and SEP
entry constants); the auxiliary fixed $L$, $\eps$, $\chi_{\rm lag}$; then $q$.
The $C^1$ field budget after $t_0$ is retained, not proved here.
The result pays its previously unpaid strip contribution but does
not close any other joint-bootstrap input. No saturation, selection,
bridge, or consolidation statement is made.
\end{document}
